\section{Introduction}
\label{sec:intro}

Equational theories in the form of rewrites ($\ell \rewritesto r$)
 have long been used in term rewriting systems.
Equality saturation engines in particular,
 which have seen a recent resurgence,
 leverage these theories to power systems in a wide
  variety of domains like program synthesis~\citep{szalinski, herbie, spores, babble, resys},
  formal verification~\citep{coq, z3, sat19, rest-darulova, sam1, sam2}, and
  optimizing compilers~\citep{eqsat, tensat, denali, wetune, deepegg, qgym, sketch-eq}.
A principal challenge in building these systems
 is writing the rewrites themselves:
 too few rewrites can lead to missed optimizations;
 too many can make implementation and maintenance difficult.
 Furthermore, even one incorrect rewrite can compromise the soundness
 of the entire system.

\textit{Theory explorers} automatically generate
 equational theories~\citep{sat19, ruler, cav21, hipster, quickspec, hipspec, isacosy}.
These tools generally follow a three-stage approach:
\begin{enumerate}
  \item
    Enumerate terms from a given grammar,
    typically in a bottom-up,
    exhaustive manner~\citep{ruler, cvc4}.
  \item
    Generate candidate rewrite rules from the enumerated terms.
    Naively, any pair of enumerated terms may be a candidate rewrite rule.
    Prior work has used techniques like finger-printing, fuzzing,
    and symbolic execution to identify ``likely sound'' candidates~\citep{ruler,sat19,cav21, bansal}.
  \item
    Using the candidates,
     select a set of rewrite rules that are both sound and useful.
    Typically,
    this is done via a process that verifies the candidates
    and removes redundant ones.
   \citet{ruler} call this process ``minimization,''
    under the assumption
    that a smaller set of rules is more likely to be effective.
\end{enumerate}

Despite recent innovations,
  theory explorers are still not widely used.
We posit that their monolithic
 implementations make them
 too inflexible.
They are designed for idealized ``one-shot'' use cases:
 the user provides a grammar, interpreter, and verifier,
 presses a button, and out comes a ruleset (set of rewrites)
 ready for use in a rewriting or equality saturation-based system.
In reality,
 tools based on equational theories
 are not developed or maintained in this manner.
Instead,
 engineers and domain experts
 build, maintain, measure, debug, and compare rulesets
 both \emph{iteratively} over time
 and \emph{incrementally} as new features and requirements are added.
In addition,
 automated theory explorers are often intended
 to replace or augment existing (handwritten) rulesets,
 but their rigid, one-shot approach
 leaves developers with little recourse
 when the output is not 100\% satisfactory.

\begin{figure}
  \centering\includegraphics[width=0.75\linewidth]{workflow.png}
  \caption{
    \textit{(Top)}
    Typical theory explorer workflow:
      the user provides a grammar, interpreter, and rule validator
      and gets a ruleset.
    Such theory explorers are rigid and opaque;
      there is no mechanism for users to intervene
      or apply domain expertise to guide inference.
    \textit{(Bottom)}
    In contrast, the \enumo DSL
      enables the user to guide theory exploration.
    Users provide the same three inputs
      as well as a short \enumo program.
    \enumo's modular and composable operators
      make it easy to implement existing inference strategies,
      add domain-specific tweaks, or even implement new strategies.
      \vspace{-0.75em}
  }\label{fig:workflow}
\end{figure}

Existing theory explorers do not scale to the
  needs of real systems.
For example,
  $x + y \rewritesto \frac{x^2 - y^2}{x-y}$
  is useful for factoring in numerical applications.
However,
  discovering it via exhaustive enumeration,
  as used in recent approaches~\citep{sat19, ruler},
  is infeasible within even a moderately sized grammar.
Instead,
  users should be able to help guide
  the theory explorer to discover such rules.

We present a new paradigm:
  theory exploration \emph{\`a la carte},
  which breaks theory explorers down
  into a set of modular operators.
Users can programmatically compose these operators
 to easily build a theory explorer suited to their needs.
We developed \enumo,
  an embedded domain-specific language (DSL)
 in which term enumeration strategies and
 rulesets are first-class values.
Simple \enumo programs can
  generate useful rulesets that
  prior work~\citep{ruler} cannot.
\enumo's abstractions also
  inspired ``fast-forwarding,''
  a new theory exploration algorithm
  that supports domains where
  equality is undecidable (e.g., real arithmetic).

We demonstrate that \enumo
  programs can synthesize
  better rulesets compared
  to state-of-the-art tools,
  while also scaling to
  much larger grammars.
In a case study inspired by Halide's large grammar~\citep{halide},
  an \enumo program synthesized a
  ruleset that derives
  90\% of Halide's handwritten rules.
Compared to prior work in theory exploration,
  fast-forwarding enabled Herbie~\citep{herbie}
  to achieve 35\% higher accuracy when
  improving floating-point rounding error (\autoref{subsec:eval-ff})
and allowed us to find alternate implementations of trigonometric
  functions in Megalibm, another floating-point synthesis tool~\citep{megalibm}.
Finally, in the domain of constructive solid geometry,
  \enumo's synthesized rules for CAD identities allow \sz~\citep{szalinski}
  to shrink benchmarks by 87\% on average,
  closely matching the 90\% reduction achieved by
  expert-written rules.

In summary, this paper makes the following contributions:
\begin{itemize}
  \item A DSL, \enumo, that offers operators for
    generating custom workloads,
    composing theory exploration strategies,
    and manipulating rulesets (\autoref{sec:slide}).
  \item A new algorithm for ``fast-forwarding'' rules to
    infer rulesets in domains where providing
    an interpreter is infeasible (\autoref{sec:lift}).
  \item An extensive evaluation showing that,
      compared to a state-of-the-art theory explorer,
      custom workloads and ruleset composition lead
      to better rulesets (\autoref{sec:eval}).
  \item A set of end-to-end case studies demonstrating that
    \enumo's synthesized rulesets are comparable to handwritten rulesets
    across a variety of domains (\autoref{sec:eval}).
  \item A case study extending \enumo with two operators that
    incorporate large language models (LLMs) into rule and workload
    synthesis, showing that the new operators compose with \enumo's
    existing pipeline and that LLM-based and guided-search rules
    are complementary
    (\autoref{sec:case-study-llm}).
  \item A detailed account of the developer experience of
    incrementally building rulesets with \enumo (\autoref{sec:devexp}).
\end{itemize}

This paper extends our earlier conference
version~\citep{enumo-oopsla23} with an expanded treatment of
derivability and proving power (\autoref{sec:slide},
\autoref{sec:eval}), a complete presentation of fast-forwarding
together with a trigonometric case study (\autoref{sec:lift},
\autoref{sec:devexp}), and a new case study extending \enumo with
operators that incorporate LLMs (\autoref{sec:case-study-llm}).

